%% =====================================================================================
%%  Executive summary (English) of the note "Voting for a consensus with veto rights", v0.1
%% -------------------------------------------------------------------------------------
%%  Build:   cd summaries/en/executive && latexmk main.tex
%%  Layout:  summaries/shared/layout.tex.  Numbers: outputs/tables/key_numbers_en.tex, written by
%%           simulations/python/key_numbers.py -- never type a number that the note computes.
%%  Keep in step with summaries/fr/executive/main.tex (same structure, same numbers).
%% =====================================================================================
\documentclass[10pt,a4paper]{article}
%% Paths that work from this folder (local build) and from the repository root (Overleaf compiles from
%% the project root whatever the main document): files are looked up under each prefix in turn.
\makeatletter
\def\input@path{{../../}{../../../}{summaries/}{./}}
\makeatother
\input{shared/layout}
\usepackage[english]{babel}
% Generated by simulations/python/key_numbers.py from outputs/tables/*.csv -- do not edit by hand
\newcommand{\KNdeadlockUnanimity}{99.6}
\newcommand{\KNdeadlockUnanimityEleven}{86}
\newcommand{\KNcontestedSchulze}{10.1}
\newcommand{\KNcontestedQVC}{5.9}
\newcommand{\KNworstoffSchulze}{85.2}
\newcommand{\KNworstoffQVC}{88.4}
\newcommand{\KNworstoffQVCquarter}{92.8}
\newcommand{\KNutilitySchulze}{97.4}
\newcommand{\KNutilityQVC}{97.7}
\newcommand{\KNutilityQVCquarter}{96.8}
\newcommand{\KNcontestedQVCquarter}{5.9}
\newcommand{\KNdeadlockUnanimityPolarised}{97.9}
\newcommand{\KNutilityPlurality}{34}
\newcommand{\KNcondorcetQVC}{94.3}
\newcommand{\KNcondorcetQVCquarter}{83.2}
\newcommand{\KNcondorcetQVCspatial}{99.5}
\newcommand{\KNfallbackQVC}{5.9}
\newcommand{\KNfallbackQVCquarter}{27}
\newcommand{\KNclonesPVC}{54}
\newcommand{\KNclonesBorda}{24}
\newcommand{\KNclonesFallback}{36}
\newcommand{\KNclonesPlurality}{98}
\newcommand{\KNmanipSchulzeSpatial}{5}
\newcommand{\KNmanipSchulzePolarised}{7}
\newcommand{\KNmanipQVCSpatial}{11}
\newcommand{\KNmanipQVCPolarised}{27}
\newcommand{\KNmanipQVCquarterSpatial}{52}
\newcommand{\KNmanipQVCquarterPolarised}{57}
\newcommand{\KNcwPolarisedTwelve}{96}


\begin{document}

\summarytitle{Voting for a consensus with veto rights}
{Executive summary of the state-of-the-art note, version 0.1}
{Fabrice Borel-Mathurin \quad$\cdot$\quad 21 October 2026 \quad$\cdot$\quad Full note, code and data:
repository \texttt{Vote\_Condorcet}}

\begin{keybox}{The question}
A group must choose one option. \textbf{Anyone may propose as many options as they wish}, \textbf{every
member judges all of them}, and \textbf{every member has a right of veto}. The procedure must reach the
\textbf{best compromise} and must \textbf{never get the group stuck}. What does research say, and which
procedure should be used?
\end{keybox}

\section{Bottom line}
\begin{enumerate}[label=\textbf{\arabic*.}]
  \item \textbf{Do not give every member a hard veto.} Unanimity breeds delay, status-quo bias and
  holdouts; with the status quo as default it deadlocks in \KNdeadlockUnanimityEleven\% of our simulated
  decisions with 11 members and \KNdeadlockUnanimity\% with 101.
  \item \textbf{Use ``Condorcet with a blocking quota'' (QVC) with a quota of one third.} One ballot; an
  option rejected by at least a third of the members is out; the majority decides among the rest.
  \item \textbf{Accept, and contain, its one price:} a blocking third can block strategically.
  Named, reasoned vetoes and at most one revision round keep that in check.
\end{enumerate}

\section{What the research says}
\begin{enumerate}[label=\textbf{\arabic*.}]
  \item \textbf{There is no perfect rule.} Every voting rule can be manipulated (Gibbard--Satterthwaite),
  and respecting the majority sometimes means overriding a large minority: designing a procedure means
  choosing a trade-off, and it should be explicit.
  \item \textbf{Unanimity is rarely optimal.} Theory and experiments agree that it raises delay and the
  power of holdouts; majority voting with a \emph{separate} veto, used only against unacceptable outcomes,
  does better (Bouton et al., 2018; Tsebelis, 2002).
  \item \textbf{Vetoes should be held by coalitions, not individuals,} as in the Council of the European
  Union's blocking minority; consensus bodies such as the IETF treat an objection as an argument to
  address, not as a block.
  \item \textbf{Open proposals require a ``clone-proof'' rule.} If anyone can add variants of an option,
  many rules reward flooding the agenda: three variants change the outcome of Borda's rule in up to
  \KNclonesBorda\% of our elections, of fallback bargaining in \KNclonesFallback\% and of the
  proportional veto core in \KNclonesPVC\%. Schulze-type Condorcet methods are immune.
  \item \textbf{Condorcet methods give the best results when preferences have structure,} and a
  Condorcet winner (an option preferred by a majority to every other) exists in at least
  \KNcwPolarisedTwelve\% of our simulated electorates, even with twelve options.
\end{enumerate}

\begin{keybox}{Recommended procedure: QVC($\theta$), with $\theta = 1/3$}
\begin{enumerate}[label=\textbf{\arabic*.},leftmargin=*]
  \item \textbf{Nominate.} Anyone proposes any number of options. The status quo is on the ballot as an
  ordinary option, never as the default.
  \item \textbf{Vote once.} Each member ranks all options (ties allowed) and marks those they find
  unacceptable; each veto is named and states a reason.
  \item \textbf{Filter.} Options vetoed by at least a third of the members are excluded. If that excludes
  them all, the least-vetoed options are kept.
  \item \textbf{Decide.} The Schulze winner (pairwise majority) among the remaining options is elected.
  Veto counts and the pairwise matrix are published.
  \item \textbf{Revise, at most once.} Only if step 3 had to fall back: reopen nominations briefly, vote
  again; the second result is final.
\end{enumerate}
\end{keybox}

\textbf{What it guarantees} (proved in the note): it \textbf{always decides}, in one round (two at most);
the status quo gets \textbf{no default advantage}; the winner is \textbf{acceptable to more than two
thirds} of the members whenever some option is; among acceptable options, \textbf{the majority decides};
and \textbf{adding variants of an option changes nothing}.

\section{The evidence}
Monte Carlo study of 21 rules, 2{,}000 simulated decisions per setting, 101 members, 8 options. Results
for a polarised group (two blocs of 55\% and 45\%), the case in which a compromise matters most:

\begin{center}\small
\begin{tabularx}{0.92\linewidth}{X r r r}
\toprule
 & \textbf{Majority only} & \textbf{QVC, $\theta = 1/3$} & \textbf{QVC, $\theta = 1/4$} \\
 & (Schulze) & (recommended) & (more consensual) \\
\midrule
Outcome rejected by at least a third of members & \KNcontestedSchulze\% & \KNcontestedQVC\% & \KNcontestedQVCquarter\% \\
Satisfaction of the least satisfied members (0--100) & \KNworstoffSchulze & \KNworstoffQVC & \KNworstoffQVCquarter \\
Total satisfaction (0--100) & \KNutilitySchulze & \KNutilityQVC & \KNutilityQVCquarter \\
Majority's preferred option elected & 100\% & \KNcondorcetQVC\% & \KNcondorcetQVCquarter\% \\
Manipulable by 10\% of members & \KNmanipSchulzePolarised\% & \KNmanipQVCPolarised\% & \KNmanipQVCquarterPolarised\% \\
\bottomrule
\end{tabularx}
\end{center}
{\footnotesize\color{muted} Satisfaction: 0 for an average option, 100 for the best; the least satisfied
member is the one at the 10th percentile. ``Rejected by a third'': \KNcontestedQVC\% is the minimum any
rule can reach, the share of decisions in which every option is rejected by a third. Unanimity with the
status quo as default, in the same setting: deadlock in \KNdeadlockUnanimityPolarised\% of decisions.\par}

\textbf{Reading.} With a blocking third, \textbf{contested outcomes are almost halved} and the least
satisfied members do better, \textbf{at no cost in total satisfaction}; the majority's choice is set
aside only when a third of the members reject it. Lowering the quota to a quarter protects minorities
more but doubles the exposure to strategic blocking.

\begin{center}
\includegraphics[width=0.8\linewidth]{fig_theta}\\[-2pt]
{\footnotesize\color{muted} What the quota $\theta$ buys and costs. Below 20\%, nearly every option is
vetoed and the procedure falls back on the least-vetoed option; above 40\%, vetoes almost never bind.
One third balances the two.\par}
\end{center}

\section{The price, and how to contain it}
Any blocking minority can also block an option it finds acceptable, to push its own favourite: no
counting rule can prevent it (Greenberg, 1979). In our simulations a coalition of 10\% of the members can
manipulate \KNmanipQVCSpatial--\KNmanipQVCPolarised\% of decisions under QVC($1/3$), against
\KNmanipSchulzeSpatial--\KNmanipSchulzePolarised\% under majority rule alone. Safeguards:
\textbf{public vetoes with stated reasons} (a chair may void an unreasoned veto); \textbf{a quota no lower
than one quarter}; \textbf{one revision round at most}; in bodies that decide often, \textbf{a budget of
vetoes across decisions}, never within one.

\section{Decisions for the organisation}
\begin{itemize}
  \item \textbf{Quota:} one third by default (with 10 members, 4 vetoes exclude an option); one quarter
  for a more consensual culture.
  \item \textbf{Status quo:} an ordinary option by default; for irreversible decisions, require the
  winner to beat it by a stated ratio (e.g.\ 2:1), as Debian does.
  \item \textbf{Revision round:} on or off, and its length.
\end{itemize}

{\footnotesize\color{muted}\textbf{Limits.} The simulations use stylised electorates; the references were
checked against publisher pages, not full texts. Next step: test the procedure on real ranked ballots
(e.g.\ Debian's published votes).\par}

\end{document}
